\documentclass[10pt,a4paper]{article}
\usepackage[utf8]{inputenc}
\usepackage[vietnamese,english]{babel}
\usepackage[top=1.2cm,bottom=1.2cm,left=1.5cm,right=1.5cm]{geometry}
\usepackage{amsmath,amssymb}
\usepackage{enumitem}
\usepackage{xcolor}
\usepackage{hyperref}
\usepackage{titlesec}

\definecolor{hcmutblue}{RGB}{0,62,126}
\definecolor{darkgray}{RGB}{50,50,50}
\definecolor{critiquecolor}{RGB}{140,25,25}

\hypersetup{
    colorlinks=true,
    linkcolor=hcmutblue,
    citecolor=hcmutblue,
    urlcolor=hcmutblue
}

\titleformat{\section}{\normalsize\bfseries\color{hcmutblue}}{\thesection.}{0.4em}{}[\titlerule]
\titlespacing*{\section}{0pt}{4pt}{2pt}

\pagestyle{empty}

\begin{document}
\selectlanguage{vietnamese}

\begin{center}
  {\footnotesize \textbf{ĐẠI HỌC BÁCH KHOA ĐHQG-HCM -- KHOA KHOA HỌC VÀ KỸ THUẬT MÁY TÍNH}}\\[0.04cm]
  {\small \textbf{CO5151: ADVANCED AGENTIC AI (HK261)} -- GIẢNG VIÊN HƯỚNG DẪN: \textbf{TS. LÊ XUÂN BÁCH}}\\[0.12cm]
  {\large\bfseries\color{hcmutblue} ANNOTATED BIBLIOGRAPHY: AGENT DESIGN PATTERNS \& ORCHESTRATION}\\[0.04cm]
  {\footnotesize \textbf{Nhóm G4}: Đặng Lâm Tùng (Trưởng nhóm) $\cdot$ Nguyễn Trung Phong $\cdot$ Vũ Việt Hưng $\vert$ \textit{Chủ đề S1-7 -- Ngày 21/10/2026}}
\end{center}

\vspace{-0.25cm}
\noindent\textit{\footnotesize \textbf{Quy ước bắt buộc}: Liệt kê đúng 4 Core Readings và 4 Extension Papers theo đề cương; mỗi bài kèm ĐÚNG MỘT CÂU nhận xét phản biện/hoài nghi học thuật (Critical Reflection), tuyệt đối không tóm tắt nội dung.}

\section{Core Readings (Các Tài Liệu Cốt Lõi)}
\begin{enumerate}[leftmargin=1.4em, itemsep=2pt, topsep=1pt]
  \item \textbf{Anthropic. (Dec 2024).} \textit{Building Effective Agents.} Anthropic Research Whitepaper.\\
  \textcolor{critiquecolor}{\textbf{Nhận xét phản biện:}} Dù phân lập chuẩn xác các khối nguyên tử thực dụng và cổ xúy cho tính tinh gọn, tài liệu đã đánh giá thấp chi phí kỹ thuật khổng lồ của việc xử lý ngoại lệ và phục hồi trạng thái khi một mắt xích công cụ trung gian gặp lỗi âm thầm trong môi trường sản xuất thực tế.

  \item \textbf{Anthropic. (Feb 2025).} \textit{How We Built Our Multi-Agent Research System.} Anthropic Engineering Blog.\\
  \textcolor{critiquecolor}{\textbf{Nhận xét phản biện:}} Mô hình phân tách Lead/Subagent tăng vọt năng lực thăm dò song song nhưng phải đánh đổi bằng sự bùng nổ chi phí token gấp 3--8 lần và nút thắt cổ chai độ trễ nghiêm trọng khi Lead Agent phải tổng hợp ngữ cảnh khổng lồ từ các tác tử phụ.

  \item \textbf{OpenAI. (2025).} \textit{A Practical Guide to Building Agents.} OpenAI Developer Documentation.\\
  \textcolor{critiquecolor}{\textbf{Nhận xét phản biện:}} Bản hướng dẫn thiên lệch về việc trao quyền tự trị tối đa cho vòng lặp Tool Calling của mô hình mà xem nhẹ rủi ro tích lũy sai số theo hàm số mũ khi số lượng bước gọi công cụ vượt quá ngưỡng kiểm soát tiền định.

  \item \textbf{LangChain Team. (2024).} \textit{LangGraph: Multi-Agent Workflows as Cyclic State Graphs.} Technical Documentation.\\
  \textcolor{critiquecolor}{\textbf{Nhận xét phản biện:}} Cơ chế Checkpointing giải quyết triệt để tính bền bỉ và khả năng phục hồi lỗi cho tác tử chạy dài, nhưng việc cưỡng bức mọi hành vi vào đồ thị trạng thái hữu hạn đã vô tình biến hệ thống thành một luồng công việc kỹ thuật truyền thống hơn là một thực thể AI tự thích ứng.
\end{enumerate}

\section{Extension Papers (Các Công Trình Mở Rộng \& Phản Biện Chuyên Sâu)}
\begin{enumerate}[leftmargin=1.4em, itemsep=2pt, topsep=1pt, start=5]
  \item \textbf{Wang, X., Wei, J., Schuurmans, D., et al. (ICLR 2023).} \textit{Self-Consistency Improves Chain of Thought Reasoning in Language Models.}\\
  \textcolor{critiquecolor}{\textbf{Nhận xét phản biện:}} Cơ chế lấy mẫu đa dạng kết hợp bầu chọn số đông thực chất chỉ lọc được nhiễu ngẫu nhiên trong quá trình giải mã chứ hoàn toàn bất lực trước các lỗi sai hệ thống xuất phát từ lỗ hổng tri thức nội tại của mô hình nền tảng.

  \item \textbf{Huang, J., Chen, X., Mishra, S., et al. (ICLR 2024).} \textit{Large Language Models Cannot Self-Correct Reasoning Yet.}\\
  \textcolor{critiquecolor}{\textbf{Nhận xét phản biện:}} Công trình bóc trần ảo tưởng về khả năng tự sửa sai nội tại của LLM, chứng minh rằng các vòng lặp phản biện thiếu tín hiệu kiểm chứng ngoại vi không những vô ích mà còn làm suy giảm độ chính xác tổng thể do thiên kiến nịnh bợ theo prompt phản biện.

  \item \textbf{Lightman, H., Kosaraju, V., Burda, Y., et al. (OpenAI, 2023).} \textit{Let's Verify Step by Step.}\\
  \textcolor{critiquecolor}{\textbf{Nhận xét phản biện:}} Dù mô hình phần thưởng theo tiến trình (PRM) cung cấp tín hiệu phản hồi từng bước vượt trội so với giám sát kết quả cuối, chi phí gán nhãn dữ liệu nhân tạo cho từng bước suy luận tạo ra một rào cản kinh tế khổng lồ khiến phương pháp này khó mở rộng ra ngoài miền toán học và lập trình.

  \item \textbf{Zelikman, E., Wu, Y., Mu, J., \& Goodman, N. (NeurIPS 2022).} \textit{STaR: Bootstrapping Reasoning With Reasoning.}\\
  \textcolor{critiquecolor}{\textbf{Nhận xét phản biện:}} Vòng lặp tự học chỉ duy trì được tính hội tụ nhờ có nhãn đáp án đúng làm màng lọc ngoại vi ở bước giữ lại suy luận, nghĩa là khả năng mở rộng của mô hình vẫn bị giới hạn tuyệt đối bởi không gian các bài toán có sẵn ground-truth xác thực.
\end{enumerate}

\vspace{-0.15cm}
\begin{center}
  \rule{0.6\textwidth}{0.4pt}\\[0.04cm]
  {\tiny \color{darkgray}CO5151 HK261 -- Nhóm G4 cam kết tự thực hiện 100\% các câu nhận xét phản biện dựa trên phân tích bài báo gốc.}
\end{center}

\end{document}
